\begin{abstract}
Multi-tenant retrieval-augmented generation (RAG) applications copy access-control state and document content into derived artifacts: ACL copies in the vector index, semantic-cache entries, and chat memory with rolling summaries. Revoking access in the identity system does not revoke these copies, and response metadata can reveal that a restricted document exists even when its content is withheld. On a synthetic five-tenant corpus, with exact-match leak judging and 48 simulated hours after each revocation, a standard deployment with a daily ACL crawl still discloses 63\% of a revoked document's secrets (leakage magnitude, LM, 0.634), and 81\% of revocations leak something. The stale index (LM 0.63), the semantic cache (0.49) and chat memory (0.37--0.51) each leak on their own, and memory leaks persist past 48\,h in 43--64\% of events. With reranker confidence exposed, an attacker infers document existence at AUC 0.983 in an undefended deployment and 0.61 in a typical one. We propose provenance-tainted revocation: every derived artifact carries the transitive set of its source documents and is served only if the user still holds all of them, combined with eager eviction and constant-shape responses. In the testbed it gives zero measured leakage over 750 events and chance-level inference (AUC 0.495), with unchanged answer quality and 360\,ms higher median latency. Ablations identify transitive taint and ACL pre-filtering as the components that matter. A check with Qwen2.5-3B shows more leakage than the stub generator (0.743 vs.\ 0.646) and the same zero under the defense.
\end{abstract}
